\chapter{Introducción}
\label{sec.ej2:introduccion}

La modulación por frecuencia (FM) es una técnica de modulación analógica en la que la frecuencia
instantánea de la señal portadora se varía en forma lineal con respecto a la amplitud de la señal
modulante. A diferencia de la modulación en amplitud (AM), la FM presenta una mayor inmunidad al
ruido y a las interferencias, a costa de un mayor ancho de banda de transmisión.

\section{Objetivos del Trabajo Práctico}
Los objetivos principales de este trabajo práctico son:
\begin{itemize}
    \item Implementar y simular un sistema completo de modulación y demodulación de FM en Multisim.
    \item Analizar el proceso de demodulación por derivación (discriminador de frecuencia).
    \item Evaluar el espectro de frecuencia de la señal FM en función del índice de modulación.
    \item Determinar matemáticamente y mediante simulación el límite máximo del índice de modulación
          para evitar distorsión en la etapa de derivación.
    \item Observar y caracterizar las formas de onda en los puntos clave del circuito (modulador,
          derivador, detector de envolvente y filtro de salida).
\end{itemize}

\section{Descripción del Sistema}
El sistema bajo estudio se divide en dos bloques principales: el modulador y el demodulador. El
esquema general del proceso a implementar se muestra a continuación:
\begin{itemize}
    \item \textbf{Modulador FM}: Genera la señal modulada a partir de una señal portadora de
          $10\kHz$ y una señal modulante sinusoidal de $1\kHz$, con un índice de modulación inicial
          $m_f = 5$.
    \item \textbf{Derivador ($d/dt$)}: Convierte las variaciones de frecuencia en variaciones de
          amplitud, generando una señal que combina modulación de amplitud y frecuencia (AM-FM).
    \item \textbf{Detección Asincrónica}: Realiza la rectificación de media onda de la señal para
          comenzar el proceso de recuperación de la envolvente.
    \item \textbf{Filtro Paso Bajo (FPB)}: Filtra las componentes de alta frecuencia (portadora y
          armónicos de rectificación), recuperando la señal de información original de $1\kHz$.
\end{itemize}
